\section{Introduction}
\label{sec:introduction}

Whether a RAG system can be trusted depends on more than the generated answer.
Retrieval relevance, evidence support, and usable provenance all matter
\citep{sun2024trustllm,ni2025trustworthyrag}.
Yet many black-box uncertainty estimators used with RAG still begin at the
answer surface.
Semantic entropy \citep{farquhar2024detecting,kuhn2023semantic},
SelfCheckGPT \citep{manakul2023selfcheckgpt}, and their RAG-integrated
variants \citep{jiang2023active,su2024dragin,zimmerman2024twotieredrag} ask,
in different forms, whether sampled answers agree.
Agreement is then used as a confidence signal, disagreement as uncertainty.
That assumption fails when retrieval repeatedly concentrates on the same
defective state. The answer is stable because the error is stable.
The problem is not that responses depend on retrieval; that is expected.
The problem is that answer agreement loses diagnostic value when the retrieval
state stops varying.
The mechanism is plain: a fixed context yields a fixed answer distribution.
The contribution is to measure its reach, namely how often silent errors
compatible with lock-in arise under deployable policies and which diagnostic
signals still expose them.

\paragraph{Retrieval-state lock-in.}
I call this failure mode \textbf{retrieval-state lock-in}: repeated samples
condition on the same defective retrieval state and therefore produce the
same wrong answer.
Two variants matter for diagnosis.
In \emph{absence} lock-in, retrieval repeatedly returns no usable graph
state, so the model answers from parametric memory.
In \emph{presence} lock-in, retrieval repeatedly returns a coherent but
wrong graph neighbourhood.
In both cases, the shared ingredient is stability, not correctness.

One presence-lock-in case asks which pharaoh the temple behind the Osireion
honours (gold: \emph{Seti~I}): the retriever anchors to \emph{Ramesses~II},
all five samples repeat that wrong name with zero answer-level entropy around a
coherent graph trace. The trace in Figure~\ref{fig:lockin_trace} shows why only
the contradicting passages raise the alarm.
Unlike ordinary noisy retrieval, where contexts vary, or the sample-consistency
view of hallucination, where confabulations diverge across samples
\citep{manakul2023selfcheckgpt,farquhar2024detecting}, the error here is stable.

Two definitions fix the scope.
\begin{definition}[Silent error]
A wrong answer with zero observed answer-state dispersion under the sampling
budget (all $N$ samples in one semantic cluster, SD-UQ at its floor).
\end{definition}
\begin{definition}[Retrieval-state lock-in]
A silent error for which the repeated samples condition on the same
\emph{defective} retrieval state, either empty (absence) or a coherent but
wrong neighbourhood (presence).
\end{definition}
Silence is a necessary symptom of lock-in, not proof of it: a silent
error can also arise from parametric overconfidence, benchmark ambiguity, or
answer-normalisation artefacts.
The analysis reports those symptom counts and uses route logs to separate the
cases that can be diagnosed from saved traces.

\begin{figure}[t]
\centering
\resizebox{\columnwidth}{!}{%
\begin{tikzpicture}[
  font=\footnotesize\sffamily,
  >={Stealth[length=2mm,width=1.5mm]},
  stage/.style={draw=black!55, line width=0.5pt, rounded corners=2pt,
                fill=white, minimum height=0.62cm, minimum width=3.45cm,
                align=center, inner xsep=5pt, inner ysep=3pt,
                font=\footnotesize\sffamily},
  bad/.style={stage, draw=badred!75!black, fill=badred!8},
  note/.style={font=\scriptsize\itshape\sffamily, text=black!55,
               align=left, anchor=west},
  raillab/.style={font=\scriptsize\sffamily, align=center, rotate=90},
  flow/.style={->, line width=0.55pt, black!60},
]
\node[stage] (q)    at (0,0)     {Question};
\node[bad]   (anc)  at (0,-1.0)  {Wrong entity anchor};
\node[bad]   (nbh)  at (0,-2.1)  {Retrieved graph\\neighbourhood};
\node[bad]   (ret)  at (0,-3.2)  {Repeated retrieval state};
\node[bad]   (ans)  at (0,-4.2)  {Repeated wrong answer};
\node[bad]   (unc)  at (0,-5.3)  {Near-zero answer\\uncertainty};

\draw[flow] (q) -- (anc);
\draw[flow] (anc) -- (nbh);
\draw[flow] (nbh) -- (ret);
\draw[flow] (ret) -- (ans)
  node[midway, right=2pt, font=\scriptsize\sffamily, text=black!50]
  {$\times\,N$ samples};
\draw[flow] (ans) -- (unc);

% ── diagnostic rail: which family observes which stage ──────────────────────
% retrieval-state family: anchor + neighbourhood
\draw[line width=2.2pt, black!30]
  (-2.00,-0.70) -- (-2.00,-2.45);
\node[raillab, text=black!45] at (-2.30,-1.58)
  {retrieval-state\\GPS $0.44$ (supported)};
% evidence-state family (the one that fires): retrieval state vs answer
\draw[line width=2.2pt, badred!75!black]
  (-2.00,-2.85) -- (-2.00,-4.55);
\node[raillab, text=badred!75!black, font=\scriptsize\bfseries\sffamily]
  at (-2.30,-3.70) {evidence-state\\SEU $=1.0$ fires};
% answer-state family (silent by construction)
\draw[line width=2.2pt, black!30]
  (-2.00,-4.95) -- (-2.00,-5.65);
\node[raillab, text=black!45] at (-2.30,-5.30)
  {answer-state\\silent};

\node[note] at (1.95,0)
  {``Osireion \ldots\ named after\\which pharaoh?'' (gold: Seti~I)};
\node[note] at (1.95,-1.0) {Ramesses~II};
\node[note] at (1.95,-2.1) {coherent, wrong builder};
\node[note] at (1.95,-3.2) {same evidence every sample};
\node[note] at (1.95,-4.2) {``Ramesses~II'' $\times 5$};
\node[note] at (1.95,-5.3)
  {$\mathrm{DSE}=0$, $\mathrm{SD\text{-}UQ}\approx 0$};
\end{tikzpicture}%
}
\caption{%
  \textbf{Worked presence-lock-in trace}, instantiated with the
  HotpotQA Osireion case also listed in
  Appendix Table~\ref{tab:lockin_gallery}.
  The retriever anchors to a strongly associated but wrong entity, the
  resulting graph neighbourhood is internally coherent, every sample receives
  the same evidence, and the repeated wrong answer presents near-zero
  answer-state uncertainty.
  The left rail shows where each diagnostic family attaches: the answer-state
  family is silent by construction, GPS remains low because the wrong answer
  is reachable in the retrieved graph, and only the evidence-state signal
  fires---the retrieved passages contradict the repeated answer.
  This traces the \emph{presence} variant; in the \emph{absence} variant
  (Section~\ref{sec:results_collapse}) the chain instead repeats an empty
  retrieval state, and the warning moves to the retrieval-state abstention
  observable.%
}
\label{fig:lockin_trace}
\end{figure}


Lock-in is easier to diagnose in knowledge-graph-augmented RAG (KG-RAG),
where a symbolic retriever can repeatedly anchor to the same entity and
traverse the same relation path.
Knowledge graphs can help on multi-hop or provenance-sensitive tasks
\citep{edge2024graphrag,gutierrez2024hipporag,hu2025grag,shen2025gear,zhu2025kg2rag,wu2025medgraphrag}.
Recent comparisons, however, find graph retrieval task-dependent and often
complementary to flat dense retrieval rather than universally superior
\citep{zhang2025ragvsgraphrag,xiang2025whengraphs,
chen2025comparingraggraphrag,peng2024graphragsurvey}.
That same symbolic structure can sharpen the failure mode, because entity-first
retrieval may keep returning the same wrong neighbourhood.

The symptom is not limited to graphs. Any system that repeatedly returns a
concentrated evidence set can produce stable wrong answers, but KG-RAG is
useful here because its retrieval state is visible: matched entities,
triples, paths, and anchors make the mechanism observable rather than merely
suspected.
The experiments use \system{},\footnote{\url{https://github.com/julka01/OntoGraphRAG}}
an ontology-guided KG-RAG framework with
entity-first routing, dense fallback, provenance-aware evidence assembly, and
graph-state logging, as the experimental setting in which to measure and
characterise lock-in.

\paragraph{Three distinct uncertainty objects.}
The central claim is that RAG confidence has three distinct objects:
\begin{enumerate}
  \item \textbf{Answer-state uncertainty} measures variation across sampled
  answers.
  \item \textbf{Evidence-state uncertainty} measures whether the retrieved
  passages locally support the generated answer.
  \item \textbf{Retrieval-state uncertainty} measures whether the retrieval
  state itself supplies graph support.
\end{enumerate}
These families are not causally independent; they are useful because each
attaches to a different object in the pipeline: the answer, the evidence, and
the retrieval state.
The triad is used throughout.\footnote{Some figures and artefact names use the
aliases \emph{output}, \emph{grounding}, and \emph{structural}.}
Each also maps to a practical trust question for deployed systems: calibrated
confidence, faithfulness to retrieved evidence, and auditable provenance.

\paragraph{Contributions.}
First, the analysis \textbf{measures the silent-error footprint compatible with
retrieval-state lock-in} under deployable policies, not only in strict stress
tests: $42\%$ of adaptive KG errors and $59\%$ of dense-retrieval errors
carry zero observed answer-state uncertainty at the operational $N{=}5$ sampling
budget, so a score using only within-question answer dispersion has no
disagreement signal to exploit: a family-level blind spot, not a stress-test
artefact.
Second, it introduces a \textbf{diagnostic decomposition} that separates
answer-state, evidence-state, and retrieval-state uncertainty.
Within the answer-state family, it also introduces SD-UQ, a low-cost
question-conditioned embedding-dispersion score that is competitive in these
low-sample RAG runs, while showing that even stronger answer-state scores cannot
diagnose lock-in alone.
The answer-state and evidence-state families are implemented directly; the
retrieval-state family is represented by a conditional graph-support
diagnostic (GPS, soft-linked and depth-matched, reported in risk orientation:
lower means stronger graph support).
Their representative scores are weakly correlated in practice (pooled Spearman
$\rho=0.03$ between answer- and evidence-state), so they carry complementary
signal rather than restating one scalar.
Third, it provides \textbf{stress-test evidence}: a strict graph-only probe
drives answer-state AUROC down to $0.233$ ($-0.52$ versus adaptive KG under
paired bootstrap), and route logs localise much of the silent error mass to
repeated empty retrieval states.
Fourth, it reports a \textbf{controlled empirical study} across five QA
families and distils the decomposition into a \textbf{conjunctive low-risk
audit rule}: treating an answer as low-risk only when all three families concur
selects a low-risk subset at $91.9\%$ pooled accuracy against a $69.7\%$ base
rate, and at $100\%$ precision under the automated judge on the clinical target
domain at $22\%$ coverage (Section~\ref{sec:results_composite}).

To separate deployment behaviour from the stress regime, the experiments compare
dense RAG, adaptive KG-RAG, and a strict graph-only KG stress test across clinical,
biomedical, and open-domain multi-hop QA snapshots.
The adaptive setting is the practical policy. The strict graph-only setting is
a mechanism probe, not a proposed retriever.
The results are mixed by design: adaptive KG retrieval is close to dense
retrieval accuracy on several snapshots (e.g., $0.946$ vs.\ $0.950$ on
RealMedQA) but does not universally improve it.
That is the condition under which the diagnostic contribution becomes precise:
when accuracy parity holds, what matters is what the retrieval layer makes
visible.

These pieces lead to the audit rule evaluated in
Section~\ref{sec:results_composite}.

\paragraph{Scope.}
Three boundaries constrain the claims below.
(1)~The absence/presence route decomposition depends on the symbolic retrieval
state that \system{} exposes; the silent-error phenomenon itself is not
graph-specific.
(2)~The strict graph-only stress test changes several pipeline components at
once, so it bounds a regime rather than isolating a single causal variable.
(3)~The fixed-subset evaluation estimates diagnostic behaviour, not lock-in
prevalence in deployed systems.
